% This is samplepaper.tex, a sample chapter demonstrating the
% LLNCS macro package for Springer Computer Science proceedings;
% Version 2.21 of 2022/01/12
%
\documentclass[runningheads]{llncs}
%
\usepackage[T1]{fontenc}
\usepackage{graphicx}
%\usepackage[round]{natbib}
%\usepackage[capitalise]{cleveref}
\usepackage{enumitem}\usepackage{paralist}
% If you use the hyperref package, please uncomment the following two lines
% to display URLs in blue roman font according to Springer's eBook style:
%\usepackage{color}
%\renewcommand\UrlFont{\color{blue}\rmfamily}
%
\begin{document}
%
\title{Efficiently explaining CSPs with Optimal Unsatisfiable Subset Optimization (Extended Abstract)}
%
%\titlerunning{Abbreviated paper title}
% If the paper title is too long for the running head, you can set
% an abbreviated paper title here
%
\author{Emilio Gamba\inst{1}\orcidID{0000-0003-1720-9428} \and
Bart Bogaerts\inst{1}\orcidID{0000-0003-3460-4251} \and
Tias Guns\inst{1,2}\orcidID{0000-0002-2156-2155}}
%
\authorrunning{E. Gamba et al.}
% First names are abbreviated in the running head.
% If there are more than two authors, 'et al.' is used.
%
\institute{Vrije Universiteit Brussel, pleinlaan 2, 1050 Brussels, Belgium\\
\email{\{emilio.gamba, bart.bogaerts, tias.guns\}@vub.be}  \and
KULeuven, Oude markt 13, 3000 Leuven, Belgium\\
\email{\{tias.guns\}@kuleuven.be}}
%
\maketitle              % typeset the header of the contribution
%
\begin{abstract}
	Recently, Bogearts et al.~\cite{bogaerts2020framework}, introduced a framework to generate a sequence of \emph{simple} explanation steps that explains all logical consequences of a Constraint Satisfaction problem (CSP). 
    An explanation in that work exploit a one-to-one correspondence between so-called \emph{non-redundant explanations} and MUSs of a derived unsatisfiable formula.
    Bogearts et al. \cite{bogaerts2020framework} identified the need for Unsatisfiable Subsets that are provably \emph{optimal} with respect to a given cost metric (that characterizes simplicity of explanations). 
	We present an implicit hitting set-based algorithm for \emph{efficiently} computing \emph{cost-optimal} \emph{constrained} unsatisfiable subsets. In the use-case of explanation for CSPs, `constrainedness', allows for the formulation structural constraints on the explanation.
%The explanation-generation algorithms presented in that work rely heavily on calls for \emph{Minimal Unsatisfiable Subsets} (MUS) \cite{marques2010minimal} of a derived unsatisfiable formula, exploiting a one-to-one correspondence between so-called \emph{non-redundant explanations} and MUSs. 
%
%
%The generation of an explanation sequence for a logic grid puzzle ranges from a few minutes to a few hours, which is more than a human would take.
%
%
%
% introduced a framework for generating a sequence of simple explanations steps for explaining solutions of constraint satisfaction problems. 
%The explanation-generation algorithms rely heavily on repeated calls for \emph{Minimal Unsatisfiable Subsets} (MUS) \cite{marques2010minimal} of a derived unsatisfiable formula, exploiting a one-to-one correspondence between so-called \emph{non-redundant explanations} and MUSs.
%
%
%
% in the framework of  
%
%
%We build on a recently proposed method .
%An explanation here is a \textit{sequence} of simple inference steps, where the simplicity of an inference step is measured by the number and types of constraints and facts used, and where the sequence explains all logical consequences of the problem. 
%The explanation-generation algorithms presented in that work rely heavily on calls for \emph{Minimal Unsatisfiable Subsets} (MUS) \cite{marques2010minimal} of a derived unsatisfiable formula, exploiting a one-to-one correspondence between so-called \emph{non-redundant explanations} and MUSs. The generation of an explanation sequence for a logic grid puzzle ranges from a few minutes to a few hours, which is more than a human would take.
%Therefore, we build on these formal foundations and tackle the main point of improvement 
%namely how to \emph{efficiently} generate explanations that are provably \emph{optimal} (with respect to the given cost metric). 
%To address this issue, we develop 
%\begin{inparaenum}[1)]
%	\item an implicit hitting set algorithm for finding \emph{optimal} \emph{constrained} unsatisfiable subsets, where `constrainedness' reduces the multiple calls for (optimal) unsatisfiable subsets to a single call;	
%	\item a method for re-using relevant information over multiple calls to these algorithms and;
%	\item methods for exploiting domain-specific information to speed-up explanation sequence generation.
%\end{inparaenum}
%We experimentally validate our algorithms on a large number of CSP instances. These experiments show that our algorithms outperform the MUS approach both in terms of quality of the explanation and in time to generate it. 
%show that this approach can be used to effectively find sequences of \textit{optimal} explanation steps for a large set of constraint satisfaction problems.
\keywords{Unsatisfiable Subset Optimization\and Explanations\and Constraint Satisfaction Problems.}
\end{abstract}
%
%
%
\section{Introduction}
	Recently \cite{bogaerts2020framework}, introduced a framework that generates a sequence of \emph{simple} explanation steps that explains all logical consequences of a Constraint Satisfaction problem (CSP). 
The explanation-generation algorithms presented in that work rely heavily on calls for \emph{Minimal Unsatisfiable Subsets} (MUS) \cite{marques2010minimal} of a derived unsatisfiable formula, exploiting a one-to-one correspondence between so-called \emph{non-redundant explanations} and MUSs. 
The generation of an explanation sequence for a logic grid puzzle ranges from a few minutes to a few hours, which is more than a human would take.
\section{Preliminaries}

\section{OCUS}

\section{Explanations}

\section{Conclusion}

\subsubsection{Acknowledgements} This research received partial funding from the Flemish Government (AI Research Program); the FWO Flanders project G070521N; and funding from the European Research Council (ERC) under the European Union’s Horizon 2020 research and innovation program (Grant No. 101002802, CHAT-Opt).

%
% ---- Bibliography ----
%
% BibTeX users should specify bibliography style 'splncs04'.
% References will then be sorted and formatted in the correct style.
%
 \bibliographystyle{splncs04}
\bibliography{krrlib,ref}

\end{document}
